\documentclass[14pt,twoside]{article}

% Paquetes a utilizar
\usepackage[utf8]{inputenc}
\usepackage[spanish]{babel}
\usepackage{amsmath,amssymb}
\usepackage{graphicx}
\usepackage{geometry}
\usepackage{fancyhdr}
\usepackage[hidelinks]{hyperref}
\usepackage{enumitem}
\usepackage{setspace}
\usepackage{parskip}
\usepackage{float}
\usepackage{booktabs}
\usepackage{tabularx}
\usepackage{array}
\usepackage{lmodern}
\usepackage{datetime2}

% Márgenes
\geometry{
  top=2.5cm,
  bottom=2.5cm,
  left=2cm,
  right=2cm
}

% Espaciado
\singlespacing

% Definir formato de fecha ISO 8601 (YYYY-MM-DD)
\newcommand{\mydate}{\DTMtoday}

% Definir título
\newcommand{\documenttitle}{Algoritmos Genéticos}

% Materia o proyecto
\newcommand{\materia}{Inteligencia Artificial y Mini-Robots (2023251)}

% Aumenta el espacio entre filas
\renewcommand{\arraystretch}{1.5}

% Encabezado y pie de página
\pagestyle{fancy}
\fancyhf{}
\fancyhead[RO]{\textbf{\materia}}
\fancyhead[LO]{\textit{\documenttitle}}
\fancyhead[RE]{\textit{\documenttitle}}
\fancyhead[LE]{\textbf{\materia}}
\fancyfoot[LE,RO]{\textit{Universidad Nacional de Colombia}}
\fancyfoot[RE,LO]{\textbf{\thepage}}
\renewcommand{\headrulewidth}{0.4pt}
\renewcommand{\footrulewidth}{0.4pt}

\begin{document}

% Título
\noindent
\begin{minipage}[c]{0.20\textwidth}
    \centering
    
\includegraphics[width=\linewidth]{0.img/UN.png}
\end{minipage}\hfill
\begin{minipage}[c]{0.75\textwidth}
    \textbf{\LARGE{\documenttitle}} \\[0.1cm]
    \hrule
    \vspace{0.1cm}
    \textbf{\materia} \\[0.3cm]
    \large{Diego Fernando Bustos Horta, \textit{dbustosh@unal.edu.co}} \\[0.1cm]
    \large{Julián Mateo Vargas Riveros, \textit{jvargasri@unal.edu.co}} \\[0.1cm]
    $\boxed{\text{\DTMtoday}}$
\end{minipage}

\vspace{0.7cm}

\section*{3.10 Ejercicios y Problemas}

El enunciado propone seis ejercicios y solicita desarrollar tres de ellos utilizando
las librerías y herramientas disponibles. En este documento se desarrollan los
Ejercicios 1, 3 y 5, buscando cubrir tres formas diferentes de aplicar un Algoritmo
Genético: optimización numérica, optimización de un problema de despacho y evolución
de cadenas de caracteres. El procedimiento sigue la estructura presentada en el
material: población inicial, evaluación de aptitud, selección, cruce, mutación y
criterio de parada.

\paragraph{Ejercicio 1. Maximización de una función numérica.}

Se desea maximizar la función
\[
f(x)=x\sin(10\pi x)+1,
\qquad x\in[0,1].
\]

Para aplicar el Algoritmo Genético Simple se utiliza una representación binaria. Se
escoge una longitud de cromosoma de $l=12$ bits, por lo que un cromosoma representa
un número decimal entre $0$ y $2^{12}-1=4095$. La decodificación utilizada es
\[
x=\frac{D}{2^{12}-1}=\frac{D}{4095},
\]
donde $D$ es el valor decimal de la cadena binaria.

Los parámetros de la simulación son:

\begin{table}[H]
\centering
\caption{Parámetros utilizados en el Algoritmo Genético del Ejercicio 1.}
\begin{tabularx}{0.82\textwidth}{>{\centering\arraybackslash}m{0.45\textwidth} >{\centering\arraybackslash}X}
\toprule
Parámetro & Valor \\
\midrule
Longitud del cromosoma, $l$ & 12 bits \\
Tamaño de población, $K$ & 50 \\
Máximo de generaciones, $M$ & 80 \\
Probabilidad de mutación, $p_m$ & 0.01 \\
Selección & Ruleta proporcional a la aptitud \\
Cruce & Un punto \\
\bottomrule
\end{tabularx}
\end{table}

La función de aptitud puede utilizarse directamente como
\[
\operatorname{Apt}(x)=f(x),
\]
porque en el intervalo considerado la función toma valores positivos. Por tanto,
los cromosomas que representan valores de $x$ con mayor valor de $f(x)$ tienen mayor
probabilidad de ser seleccionados como padres.

Después de generar la población inicial se realizan repetidamente los siguientes
pasos: decodificación de los cromosomas, cálculo de la aptitud, selección por ruleta,
cruce en un punto y mutación de bits con probabilidad $p_m$. En cada generación se
conserva para el análisis la mejor aptitud encontrada hasta ese momento.

El Algoritmo Genético obtuvo como mejor cromosoma
\[
\boxed{110110011110}
\]
con valor decimal
\[
D=3486.
\]
Por tanto,
\[
x=\frac{3486}{4095}=0.851282,
\]
y la aptitud correspondiente es
\[
f(0.851282)=1.850592.
\]
La mejor solución fue encontrada en la generación 43.

\begin{figure}[H]
\centering
\includegraphics[width=0.82\textwidth]{figuras/ejercicio1_funcion.png}
\caption{Función objetivo y solución encontrada por el Algoritmo Genético.}
\end{figure}

\begin{figure}[H]
\centering
\includegraphics[width=0.82\textwidth]{figuras/ejercicio1_convergencia.png}
\caption{Evolución de la mejor aptitud y de la aptitud promedio en el Ejercicio 1.}
\end{figure}

El resultado muestra que el Algoritmo Genético encuentra una solución próxima al
máximo de la función aun cuando el cromosoma solamente dispone de 12 bits. La
resolución de la codificación es
\[
\Delta x=\frac{1}{4095}\approx2.44\times10^{-4},
\]
por lo que existe una diferencia entre la solución discreta representada por el
cromosoma y cualquier valor real del intervalo.

\paragraph{Ejercicio 3. Despacho de energía eléctrica.}

Se tiene una empresa proveedora de energía que dispone de cuatro plantas de
generación para atender las demandas de Cali, Bogotá, Medellín y Barranquilla.
Las capacidades máximas y las demandas indicadas en el enunciado son:

\begin{table}[H]
\centering
\caption{Capacidad de generación y demanda diaria.}
\begin{tabularx}{0.82\textwidth}{>{\centering\arraybackslash}m{0.30\textwidth}>{\centering\arraybackslash}m{0.24\textwidth}>{\centering\arraybackslash}X}
\toprule
Planta / ciudad & Capacidad de generación (GW) & Demanda (GW) \\
\midrule
Cali & 3 & 4 \\
Bogotá & 6 & 3 \\
Medellín & 5 & 5 \\
Barranquilla & 4 & 3 \\
\bottomrule
\end{tabularx}
\end{table}

La demanda total es
\[
4+3+5+3=15\ \text{GW},
\]
mientras que la capacidad total disponible es
\[
3+6+5+4=18\ \text{GW}.
\]
Por tanto, existe capacidad suficiente para satisfacer la demanda.

Para representar una solución se utilizan 15 genes enteros. Cada gen representa
una unidad de $1$ GW de demanda y su valor indica la planta que suministra dicha
unidad. Los cuatro primeros genes corresponden a las cuatro unidades requeridas por
Cali, los tres siguientes a Bogotá, los cinco siguientes a Medellín y los tres
últimos a Barranquilla. De esta manera, las demandas se satisfacen exactamente por
la construcción del cromosoma y únicamente es necesario controlar las capacidades de
las plantas.

Sea $X=[x_{ij}]$ la matriz de despacho, donde $x_{ij}$ es la cantidad de energía
enviada desde la planta $i$ hacia la ciudad $j$. El costo total se calcula, siguiendo
la escala numérica de las tablas del ejercicio, como
\[
C(X)=\sum_{i=1}^{4}\sum_{j=1}^{4}
\left(g_i+t_{ij}\right)x_{ij},
\]
donde $g_i$ representa el costo de generación de la planta $i$ y $t_{ij}$ el costo
de transporte hacia la ciudad $j$.

\textit{Nota:} el enunciado presenta los costos de generación en \$/kWh y los de
transporte por GW. Para no introducir una conversión no especificada en el problema,
se conserva directamente la escala numérica de los valores proporcionados y se
interpreta $C(X)$ como un costo compuesto de comparación entre despachos.

Las restricciones de capacidad son
\[
\sum_{j=1}^{4}x_{ij}\leq \operatorname{Cap}_i,
\qquad i=1,\ldots,4.
\]

Para incorporar estas restricciones a la función de aptitud se utiliza una
penalización. Sea
\[
E(X)=\sum_{i=1}^{4}
\max\left(0,\sum_{j=1}^{4}x_{ij}-\operatorname{Cap}_i\right),
\]
la cantidad total de GW que exceden las capacidades. La aptitud se define como
\[
\operatorname{Apt}(X)=
\frac{1}{1+C(X)+10000E(X)}.
\]
Así, una solución factible, con $E(X)=0$, es favorecida frente a una solución que
viole las restricciones.

Los parámetros utilizados fueron:

\begin{table}[H]
\centering
\caption{Parámetros utilizados en el Algoritmo Genético del Ejercicio 3.}
\begin{tabularx}{0.82\textwidth}{>{\centering\arraybackslash}m{0.45\textwidth}>{\centering\arraybackslash}X}
\toprule
Parámetro & Valor \\
\midrule
Número de genes & 15 \\
Tamaño de población, $K$ & 200 \\
Máximo de generaciones, $M$ & 1000 \\
Probabilidad de mutación, $p_m$ & 0.02 \\
Selección & Ruleta proporcional a la aptitud \\
Cruce & Un punto \\
Elitismo & 2 individuos \\
\bottomrule
\end{tabularx}
\end{table}

El mejor despacho encontrado fue
\[
X^*=\begin{bmatrix}
3&0&0&0\\
1&3&0&2\\
0&0&5&0\\
0&0&0&1
\end{bmatrix}.
\]
Las filas corresponden a las plantas de Cali, Bogotá, Medellín y Barranquilla,
y las columnas a las ciudades de Cali, Bogotá, Medellín y Barranquilla.

En forma tabular:

\begin{table}[H]
\centering
\caption{Despacho de energía obtenido mediante el AG, en GW.}
\begin{tabularx}{0.92\textwidth}{>{\centering\arraybackslash}m{0.20\textwidth}>{\centering\arraybackslash}m{0.14\textwidth}>{\centering\arraybackslash}m{0.14\textwidth}>{\centering\arraybackslash}m{0.16\textwidth}>{\centering\arraybackslash}m{0.18\textwidth}}
\toprule
Planta & Cali & Bogotá & Medellín & Barranquilla \\
\midrule
Planta C & 3 & 0 & 0 & 0 \\
Planta B & 1 & 3 & 0 & 2 \\
Planta M & 0 & 0 & 5 & 0 \\
Planta Ba & 0 & 0 & 0 & 1 \\
\midrule
Demanda atendida & 4 & 3 & 5 & 3 \\
\bottomrule
\end{tabularx}
\end{table}

La solución utiliza 3 GW de la planta C, 6 GW de la planta B, 5 GW de la planta M
y 1 GW de la planta Ba. Ninguna planta supera su capacidad máxima.

El costo de generación es
\[
C_{\mathrm{gen}}=(3)(680)+(6)(720)+(5)(660)+(1)(750)
=10410.
\]
El costo de transporte es
\[
\begin{aligned}
C_{\mathrm{trans}}={}&(3)(1)+(1)(4)+(3)(1)+(2)(5)\\
&+(5)(1)+(1)(1)=26.
\end{aligned}
\]
Por tanto,
\[
\boxed{C_{\mathrm{total}}=10410+26=10436}.
\]

El AG encontró esta solución en la generación 306. La evolución del mejor costo se
muestra en la siguiente figura.

\begin{figure}[H]
\centering
\includegraphics[width=0.82\textwidth]{figuras/ejercicio3_convergencia.png}
\caption{Convergencia del mejor costo acumulado en el Ejercicio 3.}
\end{figure}

La distribución también puede visualizarse como una matriz de despacho:

\begin{figure}[H]
\centering
\includegraphics[width=0.76\textwidth]{figuras/ejercicio3_despacho.png}
\caption{Representación de la matriz de despacho obtenida.}
\end{figure}

El ejercicio muestra cómo un problema de optimización con variables discretas y
restricciones puede representarse mediante un cromosoma y resolverse evolucionando
una población de alternativas. En este caso, la penalización permite incorporar las
restricciones de capacidad dentro del proceso de evaluación.

\paragraph{Ejercicio 5. Evolución de una palabra mediante un Algoritmo Genético.}

El objetivo es generar inicialmente una población de cadenas y hacerlas evolucionar
hasta alcanzar una palabra objetivo. Para mostrar el funcionamiento del AG se toma
como palabra objetivo
\[
\boxed{\texttt{ALGORITMOS}}.
\]

Cada cromosoma tiene 10 genes y cada gen es una letra mayúscula del alfabeto inglés.
A diferencia del Algoritmo Genético Simple binario, en este caso los genes no son
bits, sino caracteres. Esto es consistente con la explicación del material, donde se
indica que un gen puede manejar diferentes tipos de datos dependiendo del problema.

La función de aptitud cuenta cuántas posiciones coinciden exactamente con la palabra
objetivo:
\[
\operatorname{Apt}(w)=
\sum_{i=1}^{10}I(w_i=\text{ALGORITMOS}_i),
\]
donde $I(\cdot)$ vale 1 cuando la condición es verdadera y 0 en caso contrario.
Por tanto,
\[
0\leq\operatorname{Apt}(w)\leq10.
\]

Los parámetros usados fueron:

\begin{table}[H]
\centering
\caption{Parámetros utilizados en el Algoritmo Genético del Ejercicio 5.}
\begin{tabularx}{0.82\textwidth}{>{\centering\arraybackslash}m{0.45\textwidth}>{\centering\arraybackslash}X}
\toprule
Parámetro & Valor \\
\midrule
Longitud del cromosoma & 10 caracteres \\
Tamaño de población, $K$ & 60 \\
Máximo de generaciones, $M$ & 100 \\
Probabilidad de mutación, $p_m$ & 0.04 \\
Selección & Ruleta proporcional a la aptitud \\
Cruce & Un punto \\
Elitismo & 2 individuos \\
\bottomrule
\end{tabularx}
\end{table}

En cada generación, las cadenas con mayor número de caracteres correctos tienen
mayor probabilidad de aportar descendencia. Durante el cruce se intercambia una
parte del contenido de dos cadenas, mientras que la mutación reemplaza una letra por
otra seleccionada aleatoriamente.

Algunas generaciones representativas fueron:

\begin{table}[H]
\centering
\caption{Evolución de la mejor cadena durante la simulación.}
\begin{tabularx}{0.86\textwidth}{>{\centering\arraybackslash}m{0.18\textwidth}>{\centering\arraybackslash}m{0.40\textwidth}>{\centering\arraybackslash}m{0.20\textwidth}X}
\toprule
Generación & Mejor cadena & Aptitud & Aptitud promedio \\
\midrule
0 & \texttt{YOGJRGCXCY} & 2/10 & 0.37 \\
10 & \texttt{AOGJRIPXOS} & 6/10 & 3.58 \\
20 & \texttt{ALGJRIPXOS} & 7/10 & 4.95 \\
30 & \texttt{ALGJRIPMOS} & 8/10 & 5.85 \\
50 & \texttt{ALGJRIPMOS} & 8/10 & 6.20 \\
70 & \texttt{ALGORISMOS} & 9/10 & 6.60 \\
80 & \texttt{ALGORISMOS} & 9/10 & 6.63 \\
84 & \texttt{ALGORITMOS} & 10/10 & 7.07 \\
\bottomrule
\end{tabularx}
\end{table}

En la generación 84 se obtuvo exactamente la palabra objetivo:
\[
\boxed{\texttt{ALGORITMOS}}.
\]
La aptitud alcanzó entonces el valor máximo de $10$.

\begin{figure}[H]
\centering
\includegraphics[width=0.82\textwidth]{figuras/ejercicio5_convergencia.png}
\caption{Evolución de la coincidencia de caracteres en el Ejercicio 5.}
\end{figure}

La simulación demuestra que el principio de evolución no está limitado a genes
binarios. En este caso, el cromosoma está formado por caracteres y los operadores
genéticos permiten transformar cadenas aleatorias hasta alcanzar la solución
objetivo.

Como elemento adicional para cumplir la parte sonora del ejercicio, la palabra final
se sintetizó mediante un sistema de texto a voz como
\texttt{ALGORITMOS}. El archivo de audio correspondiente se entrega junto con este
documento.

\paragraph{Conclusión.}

Los tres ejercicios desarrollados muestran tres usos diferentes de los Algoritmos
Genéticos. El primer ejercicio utiliza una codificación binaria para maximizar una
función numérica. El segundo representa un problema de asignación y despacho mediante
un cromosoma discreto y una función de aptitud con penalización por violaciones de
restricciones. El tercero utiliza caracteres como genes para evolucionar una palabra
hacia un objetivo determinado.

En todos los casos se mantiene la misma estructura fundamental: una población de
soluciones candidatas es evaluada mediante una función de aptitud y posteriormente
transformada mediante selección, cruce y mutación durante varias generaciones. De
esta manera, el algoritmo explora el espacio de búsqueda y utiliza la información de
las soluciones con mayor aptitud para construir nuevas soluciones.

\end{document}
